% Artigo — Template SBC (Sociedade Brasileira de Computação)
\documentclass[12pt]{article}
\usepackage[utf8]{inputenc}
\usepackage[T1]{fontenc}
\usepackage[brazil]{babel}
\usepackage{sbc-template}
\usepackage{graphicx}
\usepackage{url}
\usepackage{booktabs}
\usepackage{siunitx}
\sisetup{group-separator={.}, group-minimum-digits=4, output-decimal-marker={,}}

\sloppy

\title{Panorama das Contratações de Serviços de Computação em Nuvem\\nas Autarquias Federais do Ministério da Educação:\\Um Estudo Exploratório Baseado em Dados do PNCP}

\author{Luiz F. P. Quirino\inst{1,2}}

\address{Faculdade de Computação -- Universidade Federal de Mato Grosso do Sul (UFMS)\\
  Campo Grande -- MS -- Brasil
  \nextinstitute
  Instituto Federal de Educação, Ciência e Tecnologia de São Paulo (IFSP)\\
  São Paulo -- SP -- Brasil
  \email{luiz.quirino@ufms.br}
}

\begin{document}

\maketitle

\begin{abstract}
This paper presents an exploratory analysis of cloud computing service procurement practices across federal agencies linked to Brazil's Ministry of Education (MEC). Using public data from the National Public Procurement Portal (PNCP) API, we identify procurement patterns, modality distribution, supplier concentration, and compliance with Law No. 14.133/2021. The dataset comprises [N] procurement processes from [N] institutions between 2022 and 2026. Results reveal [main findings]. The analysis provides empirical evidence to support preliminary technical studies and contributes to the literature on digital governance in Brazilian public administration.
\end{abstract}

\begin{resumo}
Este artigo apresenta uma análise exploratória das práticas de contratação de serviços de computação em nuvem pelas autarquias federais vinculadas ao Ministério da Educação (MEC). Utilizando dados públicos da API do Portal Nacional de Contratações Públicas (PNCP), identificamos padrões de contratação, distribuição por modalidade, concentração de fornecedores e aderência à Lei nº 14.133/2021. O dataset compreende [N] processos de [N] instituições entre 2022 e 2026. Os resultados revelam [principais achados]. A análise fornece evidências empíricas para subsidiar estudos técnicos preliminares e contribui para a literatura sobre governança digital na administração pública brasileira.
\end{resumo}

\section{Introdução}
\label{sec:intro}

% TODO: Adaptar da introdução existente no relatório IFSP, com rigor acadêmico.
% Pontos: contexto (nuvem no setor público), problema (sem panorama consolidado),
% lacuna, objetivo, contribuição, organização.

\section{Trabalhos Relacionados}
\label{sec:relacionados}

% TODO: Nogueira 2021, Barros 2019, Rocha 2022, Gonçalves et al 2026, OECD 2021.
% Diferencial: foco em nuvem + recorte MEC + dados PNCP + análise de concentração.

\section{Metodologia}
\label{sec:metodologia}

% TODO: Tipo de pesquisa, fonte (API PNCP), protocolo de coleta,
% variáveis, classificação IaaS/PaaS/SaaS, técnicas de análise, limitações.
% Referência ao repositório público para reprodutibilidade.

\section{Resultados}
\label{sec:resultados}

% TODO: Subseções:
% 4.1 Visão geral (volume, evolução temporal)
% 4.2 Distribuição por modalidade
% 4.3 Tipo de serviço (IaaS/PaaS/SaaS)
% 4.4 Concentração de fornecedores (HHI)
% 4.5 Análise de valores (estimado vs homologado)
% 4.6 Maturidade por instituição

\section{Discussão}
\label{sec:discussao}

% TODO: Interpretação dos resultados, implicações para política pública,
% comparação com REGRASP (mesmo domínio, outra perspectiva),
% limitações, ameaças à validade.

\section{Conclusão}
\label{sec:conclusao}

% TODO: Síntese, contribuição prática (subsídio para ETPs),
% contribuição acadêmica (primeiro panorama empírico),
% trabalhos futuros (expansão para outros ministérios, análise textual dos TRs).

\bibliographystyle{sbc}
\bibliography{references}

\end{document}
